\section{刚体力学 III}

\subsection{快速重陀螺的旋进}

\subsubsection{旋进现象}

当陀螺旋转时受到另一个力矩，其旋转轴也会绕某个轴开始旋转，这称为\textbf{旋进现象}。

现象出现的原因是外力矩使得角动量做类似圆周运动的变化。

\subsection{作业}

\begin{problem}
    习题 7.14
    \begin{proof}
        计算刚体的转动惯量：
        $$
        I = \frac{1}{2} \cdot \frac{M l_1 R_1^2}{l_1 R_1^2 + l_2 R_2^2} R_1^2 + \frac{1}{2} \cdot \frac{M l_2 R_2^2}{l_1 R_1^2 + l_2 R_2^2} R_2^2 = \frac{M \ab(l_1 R_1^4 + l_2 R_2^4)}{2 \ab(l_1 R_1^2 + l_2 R_2^2)}
        $$
        设角加速度为 $\beta$，假设以 $m_1$ 下落为正方向，那么有方程：
        $$
        \begin{dcases}
            I \beta = F_1 R_1 - F_2 R_2 \\
            m_1 g - F_1 = m_1 \beta R_1 \\
            F_2 - m_2 g = m_2 \beta R_2\\
        \end{dcases}
        $$
        解得 $\beta = \frac{(m_1 R_1 - m_2 R_2) g}{m_1 R_1^2 + m_2 R_2^2 + \frac{M \ab(l_1 R_1^4 + l_2 R_2^4)}{2 \ab(l_1 R_1^2 + l_2 R_2^2)}}$。
    \end{proof}
\end{problem}

\begin{problem}
    习题 7.17
    \begin{solution}
        计算棒绕 $O$ 点的转动惯量：
        $$
        I = \frac{1}{12} m l^2 + \frac{1}{4} m l^2 = \frac{1}{3} m l^2
        $$
        子弹动量的变化量等于棒受到的力矩，因此：
        $$
        I \omega = \frac{1}{2} m v \cdot l \implies \omega = \frac{3 v}{2 l}
        $$
    \end{solution}
\end{problem}

\begin{problem}
    习题 7.18
    \begin{solution}
        圆盘的转动惯量：
        $$
        I = \frac{1}{2} M R^2
        $$
        设人绕轴的角速度为 $\omega_1$，圆盘绕轴的角速度为 $\omega_2$，那么根据角动量守恒：
        $$
        \begin{dcases}
            m \omega_1 R^2 + I \omega_2 = 0 \\
            (\omega_1 - \omega_2) R = v
        \end{dcases}
        $$
        解得 $\omega_1 = \frac{M v}{(2m + M) R},\ \omega_2 = - \frac{2 m v}{(2m + M) R}$。
    \end{solution}
\end{problem}

\begin{problem}
    习题 7.23
    \begin{solution}
        定性分析，碰撞后小球沿同直线做匀速直线运动，棒质心也沿这一直线做匀速直线运动，同时棒相对于质心匀速转动。

        棒绕质心的转动惯量为
        $$
        I = \frac{1}{12} M l^2
        $$
        设碰撞后小球速度为 $v_1$，棒质心速度为 $v_2$，棒绕质心角速度为 $\omega$，那么：
        $$
        \begin{dcases}
            m v_0 = m v_1 + M v_2 \\
            \frac{1}{2} m v_0 l = \frac{1}{2} m v_1 l + I \omega \\
            \frac{1}{2} m v_0^2 = \frac{1}{2} m v_1^2 + \ab(\frac{1}{2} M v_2^2 + \frac{1}{2} I \omega^2)
        \end{dcases}
        $$
        解得 $v_1 = \frac{4m - M}{4m + M} v_0,\ v_2 = \frac{2m}{4m + M} v_0,\ \omega = \frac{12m v_0}{(4m + M) l}$。
    \end{solution}
\end{problem}

\begin{problem}
    习题 7.26
    \begin{proof}
        设线轴所受摩擦力为 $f$（向左为正方向），质心加速度为 $a$（向右为正方向），线轴绕质心角加速度为 $\beta$（顺时针为正方向），线轴转动惯量为 $I$，那么有方程：
        $$
        \begin{dcases}
            a = R \beta \\
            f R - F r = I \beta \\
            F \cos \theta - f = m a \\
        \end{dcases}
        $$
        解得 $\beta = \frac{(R \cos \theta - r) F}{I + m R^2}$。
        
        那么显然，$\theta < \arccos \frac{r}{R}$ 时，线轴顺时针转动，向前滚动；$\theta > \arccos \frac{r}{R}$ 时，线轴逆时针转动，向后滚动。
    \end{proof}
\end{problem}